\begin{afterword}
\summaryhead

The post imagines a country of rationalists attacked by fanatical Barbarians, and answers the view that the rationalists must lose because each of them would rather someone else did the fighting. War is not fun, the author says, but if rationalists always lost wars, rationality could not be advocated for everyone.

In principle, the author argues, agents who reason as the author does about Newcomb's problem and the Prisoner's Dilemma coordinate on outcomes that are better for all. A community of AIs could run a lottery to pick soldiers and change their own code in advance so that the chosen will fight. For humans, the one-shot dilemma is rare: reputation, concern for others and incentives change it. A culture that does not teach defection would, the author thinks, produce volunteers; failing that, people could agree before the lottery to give selectees courage drugs and to shoot those who run. The author does not endorse present-day drafts, which the post calls a tool of kings. Soldiers must also obey unified orders, because an uncoordinated mob gets slaughtered. The purpose, the post says, is self-image: to see a society of people like oneself as able to fight and win.

\responsehead

The post answers a real objection, and its practical answer is the standard one. Its argument in principle depends on a disputed theory that the post does not give.

The argument in principle runs from Newcomb's problem through voting and the Prisoner's Dilemma to the claim that agents with the author's decision theory ``will end up at Pareto optima instead of Nash equilibria'': at outcomes that cannot be improved for anyone without making someone worse off, instead of outcomes where each does best given what the others do. Two steps have support. The post's view of voting has defenders in philosophy. Cooperation between programs that can read each other's code was a known equilibrium result before the post, for identical code, and a 2014 paper that Yudkowsky co-wrote extended it to code that differs. The general claim, and the parenthesis that common knowledge of rationality would be enough, take one side of a live dispute. The \textsc{Stanford Encyclopedia} calls the argument for cooperating with a like-minded partner ``controversial''; the other side holds that one's own choice cannot cause the other's. The post's reason is its position on Newcomb's problem, and the theory behind that position was unpublished at the time.

The practical case is better supported. The objection the post answers is the free-rider problem: each person gains by letting others do the fighting. The post's fallback is the remedy Olson named for it, ``coercion or some other special device'': a lottery, agreed on in advance and enforced by punishment. Hobbes had drawn a similar line between the soldier who has enlisted and the one who has not. The claim that people would volunteer is hedged, and history fits both it and the fallback: Britain raised its first wartime armies from volunteers and turned to conscription in 1916.

Punishment and lotteries do not separate the post's scheme from real drafts. Real armies have shot soldiers who ran, and the Vietnam draft that the post mentions was run by lottery from December 1969. What the post adds is the prior agreement of the people drafted. It states without evidence that present-day drafts are not ``collective attempts by a populace'' of this kind.

The post is open about its purpose, which is self-image, and about its limits. It ends by saying that how an army of rationalists would work is still an open question.

In short: an argument in principle that takes one side of a live dispute and depends on the author's decision theory, unpublished at the time, and a practical fallback, a lottery enforced by punishment, that is the standard remedy for free riding.
\end{afterword}
